\section{Formulace úlohy}

Cílem projektu je odhadnout popularitu hry na Steamu z informací dostupných už v době jejího
vydání: z metadat a z obrázků, kterými se hra prezentuje na své stránce (header a tři screenshoty).
Hlavní otázkou je, zda obrázky nesou informaci o popularitě navíc oproti samotným metadatům.

Steam nezveřejňuje počty prodaných kopií, proto jsem jako míru popularity zvolil celkový počet
recenzí. Jeho rozdělení je extrémně zešikmené (většina her má jednotky recenzí, několik her
miliony), a proto predikuji jeho logaritmus:
\begin{equation}
    y_i = \log\left(1 + r_i^{+} + r_i^{-}\right),
\end{equation}
kde $r_i^{+}$ a $r_i^{-}$ jsou počty pozitivních a negativních recenzí hry $i$. Chyba na této
škále vyjadřuje relativní rozdíl, takže malá i velká hra mají při trénování srovnatelnou váhu.

Úloha je regrese. Model s parametry $\theta$ předpovídá
$\hat{y}_i = f_\theta(\mathbf{m}_i, \mathbf{I}_i) \in \mathbb{R}$, kde $\mathbf{m}_i \in \mathbb{R}^{77}$
jsou metadata a $\mathbf{I}_i$ čtyři obrázky o rozměru $3 \times 128 \times 128$. Výstupem je
jedno reálné číslo bez výstupní aktivace. Parametry učím minimalizací průměrné ztráty
na trénovací množině o $n$ hrách:
\begin{equation}
    \min_\theta \; L(\theta) = \frac{1}{n} \sum_{i=1}^{n} \ell_\delta\left(\hat{y}_i, y_i\right), \qquad
    \ell_\delta(\hat{y}, y) =
    \begin{cases}
        \frac{1}{2}(\hat{y} - y)^2 & \text{pro } |\hat{y} - y| \le \delta, \\
        \delta \left(|\hat{y} - y| - \frac{1}{2}\delta\right) & \text{jinak,}
    \end{cases}
\end{equation}
kde $\ell_\delta$ je Huberova ztráta~\cite{huber1964} s $\delta = 1$. Zvolil jsem ji, protože se
u malých chyb chová jako kvadratická ztráta, ale velké chyby penalizuje jen lineárně. Hry
s nečekaně vysokou popularitou, kterou ze vstupů nelze předpovědět, tak neovládnou trénování.

Modely vyhodnocuji pomocí MAE a RMSE na logaritmické škále. MAE je hlavní metrikou, protože
se dobře interpretuje: chyba 1 znamená odhad počtu recenzí, který se od skutečnosti liší
přibližně $e \approx 2{,}7$krát. RMSE silněji penalizuje velké chyby, a ukazuje tak, jak si model
vede u extrémních případů. Model i nejlepší epochu vybírám jen podle validační množiny.